% sec_introduction.tex -- Introduction
\section{Introduction}\label{sec:introduction}

The Cayley transform
\[
\iota(x)=\frac{1+\ii x}{1-\ii x}
\]
identifies the one-point compactification of \(\RR\) with \(\TT\), and pulls
normalized Haar measure back to the Cauchy law
\[
\omega(\dd x)=\frac{1}{\pi(1+x^{2})}\dd x=P_1(x)\dd x .
\]
This is the endpoint-normalized form of the classical Cauchy/wrapped-Cauchy
M{\"o}bius correspondence in circular statistics
\cite{KentTyler1988WrappedCauchy,McCullagh1992CauchyModels,
McCullagh1996CauchyMobius,
KatoJones2010MobiusCircle,KatoMcCullagh2020SphereCauchyMobius}.
This normalization makes the fixed-law Poisson comparison
\begin{equation}\label{eq:intro-main-functional}
D_{\rm KL}\bigl(P_t*\nu\,\|\,P_t(\cdot-\bar\gamma)\bigr)
\end{equation}
the object whose large-\(t\) behaviour is studied in Cayley--Chebyshev
modes.  The input law \(\nu\) is fixed, \(\bar\gamma\) is its mean, and
\(t\to\infty\) is the Poisson scale.

The basic calculation is a quotient calculation.  With
\(x=\bar\gamma+ty\) and \(X_c:=\gamma-\bar\gamma\), translation by \(X_c\)
is encoded by
\[
R_\varepsilon(y)=\frac{1+y^2}{1+(y-\varepsilon)^2},
\qquad \varepsilon=X_c/t .
\]
The Cayley angle \(y=\tan\theta\) turns the Taylor coefficients of
\(R_\varepsilon\) into finite trigonometric, equivalently Laurent, modes.
After the matched Cauchy profile has been divided out, the quotient expansion
is therefore reduced to finite mode algebra; the KL expansion is then the
ordinary Taylor expansion of \((1+u)\log(1+u)-u\) in
\(L^\infty\cap L^1(\omega)\).

\subsection*{Main contribution and mechanism}

The new part has four pieces.  First, finite centred \((2N-2)\)th absolute
moment gives the one-dimensional coefficient expansion through order
\(t^{-2N}\), with each coefficient \(A_{2m}(\nu)\) computed by Laurent
constant-term extraction from the Cayley modes.  Second, the leading
coefficient has a sharp analytic threshold: finite centred variance is
equivalent to a finite \(t^{-4}\) matched-Poisson KL coefficient, and the
value is \(\mu_2^2/8\).  Third, in \(\RR^d\) the abstract quotient jet
\((\mathrm{MPJ}_{d,L})\) gives the tensor entropy jet, and finite centred
\(L\)-absolute moment with integer \(L\ge d+1\) is proved to imply that jet.
Fourth, under finite centred second moment, the same matched quotient
satisfies an exact Cauchy-generator Bregman dissipation identity whose tail
integral recovers the KL value.

The one-dimensional coefficient theorem, sharp threshold, multidimensional
jet, and dissipation identity are the main formal statements below.  The
finite-moment quotient remainder, explicit power-tail replacement quotient,
and multidimensional moment jet are stated separately so that the analytic
input and the algebraic entropy transfer remain visible.

\subsection*{Formal theorem statements}

The formal statements below give the exact assumptions and coefficients used
later.  The multidimensional coefficients are denoted
\(\mathcal A_m^{(d)}(\nu)\); the one-dimensional coefficients use the
unsuperscripted symbols \(A_k(\nu)\).  Throughout these statements,
\(\bar\gamma\) denotes the matched barycentre and
\(X_c:=\gamma-\bar\gamma\) denotes the centred variable.

\begin{theorem}[One-dimensional coefficient identification at finite \((2N-2)\)-moment level]%
\label{thm:intro-one-dimensional-universal-expansion}
Let \(N\ge2\), let \(\nu\) be a probability measure on \(\RR\) with mean
\(\bar\gamma\), and put \(X_c:=\gamma-\bar\gamma\).  Assume that \(X_c\)
has finite \((2N-2)\)th absolute moment, and put
\[
h_t=P_t*\nu,\qquad g_t(x)=P_t(x-\bar\gamma),\qquad
\mu_n=\int X_c^n\,\dd\nu(\gamma).
\]
Then
\begin{equation}\label{eq:intro-one-dimensional-transfer}
D_{\mathrm{KL}}\bigl(P_t*\nu\,\|\,P_t(\cdot-\bar\gamma)\bigr)
=
\sum_{m=2}^{N}A_{2m}(\nu)t^{-2m}+o(t^{-2N}).
\end{equation}
Each \(A_{2m}\) is the finite-moment evaluation of a universal polynomial in
the centred moments, computed by the Cayley mode integrals in
Theorem~\ref{thm:universal-entropy-coefficients}; all retained odd
coefficients vanish.  More explicitly, for \(2\le m\le N\),
\[
A_{2m}(\nu)=
\sum_{r=2}^{m}
\frac{(-1)^r}{r(r-1)}
\sum_{\substack{n_1+\cdots+n_r=2m\\2\le n_j\le 2N-2}}
\mu_{n_1}\cdots\mu_{n_r}
\int_{\RR}u_{n_1}\cdots u_{n_r}\,\dd\omega .
\]
\end{theorem}

\begin{theorem}[Sharp leading tail threshold in dimension one]%
\label{thm:intro-sharp-leading-threshold}
Let \(\nu\) be a probability measure on \(\RR\) with finite mean
\(\bar\gamma\), and put
\[
h_t=P_t*\nu,\qquad g_t(x)=P_t(x-\bar\gamma),\qquad
H(t)=D_{\mathrm{KL}}(h_t\|g_t).
\]
The last quantity is the extended KL functional of
\eqref{eq:extended-kl-definition}, so \(H(t)\in[0,\infty]\) for every
\(t>0\).
If the centred variance \(\mu_2=\int(\gamma-\bar\gamma)^2\,\dd\nu(\gamma)\)
is finite, then \(H(t)<\infty\) for every \(t>0\) and
\[
\lim_{t\to\infty}t^4H(t)=\frac{\mu_2^2}{8}.
\]
If the centred variance is infinite, then
\[
\lim_{t\to\infty}t^4H(t)=+\infty.
\]
The infinite-variance assertion is an identity in the extended interval:
it includes the possibility that \(H(t)=+\infty\) at some, or all,
large times.
Equivalently, finite centred second moment is necessary and sufficient for a
finite leading \(t^{-4}\) matched-Poisson KL coefficient, and the finite value
is the universal coefficient \(A_4(\nu)=\mu_2^2/8\).
\end{theorem}

Theorems~\ref{thm:intro-one-dimensional-universal-expansion}
and~\ref{thm:intro-sharp-leading-threshold} are the finite-moment coefficient
theorem and the sharp finite-variance threshold, respectively.  Tail
replacements are intentionally isolated in
Theorem~\ref{thm:tail-replacements-proved-here}.  For the explicitly
displayed symmetric power-tail density
\[
f(x)=g(x)+c|x|^{-p-1}{\bf1}_{|x|\ge R},
\]
and for perturbations satisfying
\eqref{eq:perturbative-power-tail-lower-moments}--\eqref{eq:perturbative-power-tail-far-remainder},
in particular any perturbation with
\(|r(x)|\le C|x|^{-p-1-\eta}\) for some \(\eta>0\),
the quotient has a top-mode truncated-moment term
\(u_p(y)t^{-p}M_p(t)\); Theorem~\ref{thm:tail-replacements-proved-here}
then transfers this to the entropy layer
\(\kappa_N\mu_2t^{-2N}M_p(t)\).

\begin{theorem}[Finite-moment multidimensional Poisson KL asymptotic]%
\label{thm:intro-multidimensional-jet}
\label{cor:intro-multidimensional-finite-moment-entropy-jet}
Let \(d\ge1\), let \(L\in\mathbb N\) satisfy \(L\ge d+1\), and let \(\nu\)
be a probability measure on \(\RR^d\) with mean \(\bar\gamma\) and finite
centred \(L\)-absolute moment.  Here
\[
P_t^{(d)}(x)=
\frac{\Gamma((d+1)/2)}{\pi^{(d+1)/2}}\,
\frac{t}{(t^2+|x|^2)^{(d+1)/2}}
\]
is the normalized radial Poisson kernel on \(\RR^d\).  Put
\[
h_t=P_t^{(d)}*\nu,\qquad
g_t=P_t^{(d)}(\cdot-\bar\gamma),\qquad
X_c:=\gamma-\bar\gamma,
\]
and write
\[
\mu_\alpha=\EE[X_c^\alpha]\qquad(|\alpha|\le L).
\]
Then
\[
D_{\rm KL}(h_t\|g_t)
=\sum_{m=4}^{L+2}\mathcal A_m^{(d)}(\nu)t^{-m}
+o(t^{-(L+2)}),
\]
where the tensor coefficients are
\[
\mathcal A_m^{(d)}(\nu)=
\sum_{r=2}^{\lfloor m/2\rfloor}
\frac{(-1)^r}{r(r-1)}
\sum_{\substack{|\alpha_1|,\ldots,|\alpha_r|\ge2\\
|\alpha_1|+\cdots+|\alpha_r|=m\\
|\alpha_j|\le L}}
\mu_{\alpha_1}\cdots\mu_{\alpha_r}
\int_{\RR^d}\prod_{j=1}^rU_{\alpha_j}^{(d)}(y)\,\Omega_d(\dd y).
\]
All odd retained coefficients vanish:
\[
\mathcal A_m^{(d)}(\nu)=0
\qquad\text{for every odd }m,\quad 4\le m\le L+2.
\]
Moreover the proof gives the explicit quotient remainder
\[
\delta_t^{(d)}(y)=
\frac{h_t(\bar\gamma+ty)}{P_t^{(d)}(ty)}-1
=
\sum_{2\le|\alpha|\le L}
\mu_\alpha U_\alpha^{(d)}(y)t^{-|\alpha|}
+\rho_{L,t}^{(d)}(y),
\]
with
\[
\|\rho_{L,t}^{(d)}\|_\infty+
\|\rho_{L,t}^{(d)}\|_{L^1(\Omega_d)}
\le C_{d,L}t^{-L}\EE\!\left[
|X_c|^L\frac{|X_c|}{t}{\bf1}_{\{|X_c|\le t\}}
{}
+{}|X_c|^L{\bf1}_{\{|X_c|>t\}}
\right]
=o(t^{-L}),
\]
and \(\|\delta_t^{(d)}\|_\infty=O(t^{-2})\).  The first term in this bound is
the local Taylor remainder; the second is the large-translation contribution,
where \(L\ge d+1\) is used.

When
\[
\Sigma=\EE[X_cX_c^{\mathsf T}],\qquad
\Sigma_0=\Sigma-d^{-1}(\operatorname{tr}\Sigma)I_d ,
\]
the leading coefficient is
\begin{equation}\label{eq:intro-multidim-leading}
\mathcal A_4^{(d)}(\nu)
=
c_d^{\rm iso}(\operatorname{tr}\Sigma)^2
+c_d^{\rm tr}\|\Sigma_0\|_{\rm HS}^2,
\end{equation}
where
\[
c_d^{\rm iso}=
\frac{3(d+1)(7d+9)}
{4d(d+3)(d+5)(d+7)},
\qquad
c_d^{\rm tr}=
\frac{3(d+1)(d+3)}
{4(d+5)(d+7)} .
\]
Finite covariance alone identifies the displayed quadratic form; the KL
asymptotic above uses the stronger quotient expansion just displayed, supplied
here by the finite centred \(L\)-moment hypothesis.
\end{theorem}

\begin{proposition}[Abstract multidimensional entropy transfer]%
\label{prop:intro-multidimensional-mpj-transfer}
Let \(d\ge1\), let \(L\in\mathbb N\) satisfy \(L\ge d+1\), and let \(\nu\)
be a probability measure on \(\RR^d\) with mean \(\bar\gamma\).  Assume that
the centred tensor moments \(\mu_\alpha=\EE[X_c^\alpha]\) exist for
\(2\le|\alpha|\le L\), and that
\[
\begin{aligned}
(\mathrm{MPJ}_{d,L}):\qquad
\delta_t^{(d)}(y)
&=
\sum_{2\le|\alpha|\le L}
\mu_\alpha U_\alpha^{(d)}(y)t^{-|\alpha|}
+\rho_{L,t}^{(d)}(y),\\
\|\rho_{L,t}^{(d)}\|_\infty+
\|\rho_{L,t}^{(d)}\|_{L^1(\Omega_d)}
&=o(t^{-L}),\qquad
\|\delta_t^{(d)}\|_\infty=O(t^{-2}),
\end{aligned}
\]
with \(\delta_t^{(d)}\) defined as in
Theorem~\ref{thm:intro-multidimensional-jet}.  Then the same entropy expansion
and coefficient formula of Theorem~\ref{thm:intro-multidimensional-jet}
follow.  The finite centred \(L\)-moment theorem above is obtained by first
proving this \((\mathrm{MPJ}_{d,L})\) transfer and then proving the displayed
large-translation remainder bound from the moment assumption.
\end{proposition}

\begin{theorem}[Poisson quotient KL tail identity]%
\label{thm:intro-kl-tail-identity}
Let \(\nu\) be a probability measure on \(\RR\) with mean \(\bar\gamma\) and
finite centred second moment.  For
\[
h_t=P_t*\nu,\qquad g_t(x)=P_t(x-\bar\gamma),\qquad u_t=h_t/g_t,
\]
define
\[
\mathcal I_\nu(t)=
\frac1\pi\iint_{\RR^2}
\frac{g_t(x)\left(
u_t(x)\log\!\bigl(u_t(x)/u_t(y)\bigr)-u_t(x)+u_t(y)
\right)}
{(x-y)^2}\dd x\dd y .
\]
Then \(\mathcal I_\nu\in L^1_{\rm loc}(0,\infty)\) and
\[
D_{\mathrm{KL}}\bigl(P_t*\nu\,\|\,P_t(\cdot-\bar\gamma)\bigr)
=\int_t^\infty \mathcal I_\nu(s)\dd s
\]
for every \(t>0\).
\end{theorem}

This is the paper's de Bruijn/Stam-type result for the matched Cauchy
scale.  The foundational Gaussian line is de Bruijn's heat-flow entropy
identity and Stam's Fisher-information entropy-power inequality
\cite{DeBruijn1953,Stam1959,Blachman1965,Barron1986,JohnsonBarron2004,
MadimanBarron2007}.  Moving-reference relative-entropy identities, where
both the numerator and reference density evolve under a compatible flow, are
represented by \cite{Verdu2010Mismatched,
HirataNemotoYoshida2012RelativeEntropy,Yoshida2017RelativeEntropyDiffusion}.
Theorem~\ref{thm:intro-kl-tail-identity} gives the Cauchy-generator
specialization for the nonstationary pair
\[
(h_t,g_t)=\bigl(P_t*\nu,P_t(\cdot-\bar\gamma)\bigr),
\]
including the direct cutoff-domain verification for the moving quotient
\(u_t=h_t/g_t\) and the tail representation that transfers the Cayley KL
coefficients to the dissipation \(\mathcal I_\nu\).

\begin{theorem}[Explicit \(N=4\) coefficient layer through order eight]%
\label{thm:intro-explicit-n4-layer}
Let \(\nu\) be a probability measure on \(\RR\) with mean \(\bar\gamma\),
finite variance \(\sigma^2\ge0\), and finite centred sixth absolute moment.
Write
\[
X_c:=\gamma-\bar\gamma,\qquad
\mu_j:=\int_{\RR}X_c^j\,\dd\nu(\gamma)\quad (j\le6),
\qquad
\sigma^2=\mu_2 .
\]
Then
\[
D_{\mathrm{KL}}\bigl(P_t*\nu\,\|\,P_t(\cdot-\bar\gamma)\bigr)
=
\frac{\sigma^4}{8t^4}
+\frac{A_6(\nu)}{t^6}
+\frac{A_8(\nu)}{t^8}
+o(t^{-8}),
\]
where
\[
A_6(\nu)
=\frac{\sigma^6+6\mu_3^2-8\sigma^2\mu_4}{64},
\]
and
\[
A_8(\nu)
=
\frac{3\sigma^8-12\sigma^4\mu_4+9\sigma^2\mu_3^2
+12\sigma^2\mu_6-30\mu_3\mu_5+20\mu_4^2}{256}.
\]
If the centred seventh absolute moment is finite, then
the last remainder improves to \(O(t^{-9})\).  When \(\sigma^2=0\),
\(\nu=\delta_{\bar\gamma}\) and the expansion is the trivial identity
\(0=0\).
\end{theorem}

\subsection*{Related work and theorem-level comparison}

The nearest recent \emph{Journal of Functional Analysis} comparators and the
present theorem groups are most directly compared by the functional being
expanded, the object that evolves, the hypotheses, and the asymptotic output.
The table below records those mathematical axes.  Appendix~\ref{app:related-work-scope}
keeps the same axes for the recent-JFA comparison and then lists older or
non-JFA background used for terminology and proof language.

\begingroup
\footnotesize
\setlength{\LTleft}{0pt}
\setlength{\LTright}{0pt}
\setlength{\tabcolsep}{1.8pt}
\renewcommand{\arraystretch}{1.14}
\begin{longtable}{p{0.17\linewidth}p{0.20\linewidth}p{0.17\linewidth}p{0.22\linewidth}p{0.16\linewidth}}
\hline
\(\text{paper}\) & \(\text{functional}\) & \(\text{evolving object}\) &
\(\text{hypotheses}\) & \(\text{asymptotic output}\)\\
\hline
\endfirsthead
\hline
\(\text{paper}\) & \(\text{functional}\) & \(\text{evolving object}\) &
\(\text{hypotheses}\) & \(\text{asymptotic output}\)\\
\hline
\endhead
Bobkov--G{\"o}tze
\cite{BobkovGotzeEsscher2025}
&
\(D_\infty\) / Esscher information distance
&
Esscher transforms of normalized sums
&
CLT-type normalization and moment/tail hypotheses
&
Information convergence and quantitative normal-approximation rates\\
\hline
Anker--Papageorgiou--Zhang
\cite{AnkerPapageorgiouZhang2023HeatSymmetricSpaces}
&
Linear heat profile and heat-kernel asymptotics
&
Heat equation on \(G/K\)
&
\(L^1\), symmetry, or geometric hypotheses appropriate to the symmetric space
&
Large-time profile expansion\\
\hline
Bogdan--Gutowski--Pietruska-Pa{\l}uba
\cite{BogdanGutowskiPietruskaPaluba2025PolarizedHardyStein}
&
Polarized Bregman/Hardy--Stein energy
&
Semigroup or Dirichlet-form data
&
Markov-semigroup and convexity hypotheses
&
Exact energy identity\\
\hline
Present paper,
Theorems~\ref{thm:intro-one-dimensional-universal-expansion}
and~\ref{thm:intro-sharp-leading-threshold}
&
\(D_{\rm KL}(P_t*\nu\|P_t(\cdot-\bar\gamma))\)
&
One fixed law; increasing Poisson scale \(t\)
&
Finite centred \((2N-2)\)-moment for retained coefficients; finite variance
exactly for a finite leading coefficient
&
Explicit Laurent/KL coefficients \(A_{2m}\) and \(A_4=\mu_2^2/8\)\\
\hline
Present paper,
Theorem~\ref{thm:tail-replacements-proved-here}
&
\(D_{\rm KL}(P_t*\nu\|P_t(\cdot-\bar\gamma))\) with a top quotient
tail mode
&
One fixed power-tail or perturbative power-tail law; increasing Poisson scale
&
Displayed symmetric power-tail density or perturbations satisfying
\eqref{eq:perturbative-power-tail-lower-moments}--\eqref{eq:perturbative-power-tail-far-remainder}
&
Explicit top-mode replacement involving \(M_p(t)\), including the stated
borderline logarithmic layers\\
\hline
Present paper,
Theorem~\ref{thm:intro-multidimensional-jet} and
Proposition~\ref{prop:intro-multidimensional-mpj-transfer}
&
Multidimensional matched radial-Poisson KL
&
One fixed law on \(\RR^d\); increasing radial-Poisson scale
&
\((\mathrm{MPJ}_{d,L})\) with \(L\in\NN\), \(L\ge d+1\); finite centred
\(L\)-absolute moment as a sufficient condition
&
Tensor KL coefficient jet and isotropic/traceless covariance coefficient\\
\hline
Present paper,
Theorem~\ref{thm:intro-kl-tail-identity}
&
Matched-Poisson KL and Cauchy Bregman dissipation \(\mathcal I_\nu(t)\)
&
Moving quotient \(u_t=(P_t*\nu)/P_t(\cdot-\bar\gamma)\)
&
Finite centred variance
&
Exact tail identity
\(D_{\rm KL}=\int_t^\infty\mathcal I_\nu(s)\,\dd s\) and coefficient transfer\\
\hline
Present paper,
Theorem~\ref{thm:intro-explicit-n4-layer} and
Corollary~\ref{cor:intro-defect-certificate}
&
Coefficient layer \(A_4,A_6,A_8\) and the defect \(\Delta_6\)
&
One fixed law with fixed variance
&
Finite centred sixth moment for \(A_8\); finite fourth moment for
\(\Delta_6\)
&
Closed coefficient formulae, \(\Delta_6\ge0\), and a finite-moment
coefficient-level \(W_2\) coupling consequence\\
\hline
\end{longtable}
\endgroup

For the first row, the coefficient conclusion is the explicit constant-term
formula
\[
A_{2m}(\nu)=
\sum_{r=2}^{m}\frac{(-1)^r}{r(r-1)}
\sum_{n_1+\cdots+n_r=2m}
\mu_{n_1}\cdots\mu_{n_r}
\int u_{n_1}\cdots u_{n_r}\,d\omega .
\]

\subsection*{Cayley coefficient mechanism}

The opening quotient calculation is made formal in
Theorem~\ref{thm:cayley-chebyshev-mode}, which gives the Cayley--Chebyshev
mode formula.  Theorem~\ref{thm:universal-entropy-coefficients} defines the
entropy coefficients by finite Cauchy-weighted mode integrals, and
Theorem~\ref{thm:finite-fourier-mode-integral-algorithm} reduces those
integrals to Laurent constant-term extraction.

The finite-moment quotient remainder is proved in
Lemma~\ref{lem:exact-moment-density-truncation}.  The quotient-to-KL step is
split into the abstract finite-jet transfer
Proposition~\ref{prop:finite-quotient-jet-implies-kl-coefficient-jet}, the
retained-order density-jet transfer
Proposition~\ref{prop:finite-density-jet-entropy-transfer}, and the resulting
finite-moment Poisson coefficient theorem
Theorem~\ref{thm:poisson-kl-even-orders}.  The finite Laurent arithmetic
needed for the explicit \(N=4\) layer is displayed in
Lemmas~\ref{lem:n4-normalized-laurent-representatives}--\ref{lem:n4-entropy-coefficient-substitution}.

\subsection*{Further consequences}

\begin{corollary}[Coefficient-defect coupling at the fourth-moment layer]%
\label{cor:intro-defect-certificate}
Let \(\nu\) have mean \(\bar\gamma\), variance \(\sigma^2>0\), and finite
centred fourth absolute moment.  Put
\[
X_c:=\gamma-\bar\gamma,\qquad
\mu_j:=\int_{\RR}X_c^j\,\dd\nu(\gamma)\quad (j\le4),
\qquad
\sigma^2=\mu_2 .
\]
The \(N=3\) case of
Theorem~\ref{thm:intro-one-dimensional-universal-expansion} identifies
\[
A_6(\nu)=\frac{\sigma^6+6\mu_3^2-8\sigma^2\mu_4}{64}
\]
as the genuine \(t^{-6}\) coefficient.  Define the sixth-order defect
\[
\Delta_6(\nu)=-\frac{7}{64}\sigma^6-A_6(\nu).
\]
Equivalently, for \(X=X_c/\sigma\), \(\beta_3=\EE X^3\), and
\(Q_2=X^2-1-\beta_3X\),
\[
\Delta_6(\nu)=
\sigma^6\left(\frac{\beta_3^2}{32}+\frac{\|Q_2\|_2^2}{8}\right).
\]
Corollary~\ref{cor:coefficient-defect-coupling} couples this fixed-variance
coefficient defect to
\(\rho_\sigma=\frac12(\delta_{-\sigma}+\delta_\sigma)\) in \(W_2\).
\end{corollary}

This corollary uses the displayed fourth-moment coefficient algebra:
the identity for \(A_6(\nu)\) is combined with a coupling of the normalized
law to the symmetric two-point law.

\begin{corollary}[Leading multidimensional KL coefficient under a finite jet hypothesis]%
\label{cor:intro-multidimensional-leading-kl-coefficient}
Let \(d\ge1\), let \(L\in\mathbb N\) satisfy \(L\ge d+1\), and assume that
\(\nu\) has mean \(\bar\gamma\), centred tensor moments
\(\mu_\alpha\) for \(2\le|\alpha|\le L\), and satisfies
\[
\begin{aligned}
(\mathrm{MPJ}_{d,L}):\qquad
\delta_t^{(d)}(y)
&=\sum_{2\le|\alpha|\le L}
\mu_\alpha U_\alpha^{(d)}(y)t^{-|\alpha|}
+\rho_{L,t}^{(d)}(y),\\
\|\rho_{L,t}^{(d)}\|_\infty+
\|\rho_{L,t}^{(d)}\|_{L^1(\Omega_d)}
&=o(t^{-L}),\qquad
\|\delta_t^{(d)}\|_\infty=O(t^{-2}).
\end{aligned}
\]
By Theorem~\ref{thm:multidimensional-density-quotient}, finite centred
\(L\)th absolute moment is a sufficient condition for this jet.
With \(A_{4,\mathrm{quad}}^{(d)}(\nu)\) as in
\eqref{eq:intro-multidim-leading},
\[
D_{\rm KL}\bigl(P_t^{(d)}*\nu\,\|\,P_t^{(d)}(\cdot-\bar\gamma)\bigr)
=A_{4,\mathrm{quad}}^{(d)}(\nu)t^{-4}+o(t^{-4}).
\]
This is the abstract-jet version of the leading coefficient statement.  The
finite-\(L\)-moment sufficient condition is already included in
Theorem~\ref{thm:intro-multidimensional-jet}.  The quotient remainder
hypothesis turns the finite-covariance quadratic form into an asymptotic
coefficient.
\end{corollary}

\subsection*{Scope and limitations}

The KL functional used throughout is the extended nonnegative
Kullback--Leibler functional of \eqref{eq:extended-kl-definition}
\cite{KullbackLeibler1951,CoverThomas2006}.  Under finite centred second
moment in dimension one the matched-Poisson KL value is finite for every
\(t>0\); without that hypothesis the sharp-threshold theorem is interpreted
in \([0,\infty]\), so the conclusion
\(t^4D_{\rm KL}(P_t*\nu\|P_t(\cdot-\bar\gamma))\to+\infty\) includes the
case where the KL value is already \(+\infty\) at some large times.  The
finite-moment hierarchy identifies universal polynomial coefficients only
when the needed centred moments exist; for the explicit symmetric power-tail
laws treated in the main text, missing top moments belong instead to the
proved logarithmic replacement layers.  Theorem~\ref{thm:tail-replacements-proved-here}
is an explicit symmetric power-tail examples theorem for the displayed
density classes; Corollary~\ref{cor:conditional-slowly-varying-borderline-tail}
adds a slowly varying borderline tail only under the stated Potter-type
quotient controls.  These statements are not a general slowly varying or
regularly varying replacement theorem.  In several dimensions, the finite-covariance
proposition is the quadratic coefficient calculation, whereas the KL
asymptotic requires the separate quotient-jet hypothesis
\((\mathrm{MPJ}_{d,L})\), for which finite centred \(L\)-moment with integer
\(L\ge d+1\) is a sufficient condition.  Corollary~\ref{cor:intro-defect-certificate}
is a finite-moment consequence of the coefficient algebra, not an
optimizer-stability theorem for a sharp functional inequality and not a
claim that a small matched-Poisson KL value alone controls \(W_2\).  The
dissipation identity is proved in Section~\ref{sec:doob-phi-entropy} by the
Cauchy cutoff Green formula.  That section also records that pointwise
dissipation asymptotics need additional differentiated-remainder input; the
coefficient theorems use only tail integrals and fixed-window averages.

\subsection*{Organization}

Section~\ref{sec:preliminaries} fixes notation.  Sections~\ref{sec:cayley-gate}
and~\ref{sec:haar-pullback} record the Cayley and Haar-pullback
normalizations.  Section~\ref{sec:entropy-asymptotics} proves the
multidimensional entropy jets, the one-dimensional coefficient hierarchy, the
sharp leading variance threshold, the explicit coefficient layer through
order eight, and the sixth-order coefficient-defect consequence.
Section~\ref{sec:doob-phi-entropy}
derives the Cauchy/KL Bregman tail identity by cutoff Green calculus,
verifies the Poisson quotient hypotheses, and then proves the dissipation-tail
coefficient expansion.
Appendix~\ref{app:related-work-scope} gives the theorem-level
related-work comparison;
Appendix~\ref{app:entropy-integrals} records the weighted trigonometric
normalization behind the Laurent arithmetic;
Appendix~\ref{app:haar-pullback-uniqueness}
records the endpoint-normalized Cayley uniqueness statement.
